% !TEX root = ../../../main.tex
% Block card: preprocessing.beads -- "BEADS baseline" (Part I, the blocks chapter).
\cardfield{What it does}
Removes slow drift from a signal by splitting it into three layers at once (a smooth baseline,
a layer of rare events, and noise) and subtracting only the baseline.

\cardfield{When to use it}
When the recording has big drops or spikes that a moving average would be dragged towards.
Imagine tracing the horizon in a photo with a few tall buildings: a ruler that averages
everything in view bends up at each building, while someone told ``buildings are rare and
the horizon is smooth'' draws the horizon straight underneath them. BEADS is that second
person. It comes from chromatography \citep{ning2014chromatogram}; here it is adapted for long,
slow fungal recordings.

\cardfield{What goes in, what comes out}
In goes a span of one channel. Out comes the same span, same length and rate, with the baseline
subtracted. The block page also reports how well the fit went and any warnings.

\cardfield{Settings that matter}
\emph{baseline\_period\_s} (two hours by default) is the most important: anything slower than
this is treated as drift and removed, anything faster is kept. Set it above the slowest event
you care about. \emph{polarity} tells it which way your events point: \emph{drop} (the default)
for troughs, \emph{spike} for peaks, \emph{both} for no preference. \emph{sparsity} says how
rare events must be: raise it if too much of the trace is treated as events.

\cardfield{Watch out for}
An event as slow as the baseline period is removed with the drift: no method can tell them
apart. On a span made mostly of events, the baseline settles partway into them, so event shapes
and depths stay trustworthy but absolute levels do not. An event close to either end of the
span can bend the baseline there; the block warns when it sees one. Read the warnings it prints.
